
\section*{Ex.1 $|x^2+1|<x+1$}
Per poder llevar el valor absolut passam:
\begin{equation*}
|x^2+1|<x+1 \ \Longleftrightarrow \ x^2+1<x+1 \ \wedge \ x^2+1>-x-1
\end{equation*}
\\
Començam per la primera part:
\\
\begin{equation*}
x^2+1<x+1 \ \Leftrightarrow \ x^2-x-2<0 
\end{equation*}
\\
Passam la ineqüació a una equació de segon grau: $x^2-x-2=0$ que té per solucions $x=2 \ \wedge \ x=-1$. \\
\\
Cream una taula pels valors:
\begin{table}[h!]
\begin{center}
\begin{tabular}{c|c|c}
$\oplus$ & $\ominus$ & $\oplus$\\
\hline
$(-\infty,-1)$ & $(-1,2)$ & $(2,+\infty)$
\end{tabular}
\end{center}
\end{table}
\\
Per tant la primera solució és $(-1,2)$.
\\
\\
Ens centram ara en la segona part:
\begin{equation*}
x^2+1>-x-1 \ \Leftrightarrow \ x^2+x>0 \ \Leftrightarrow \ x(x+1)>0
\end{equation*}
Ara resolem la casuítica:
\begin{equation*}
\left.
\begin{array}{cc}
\left\lbrace
\begin{array}{cc}
x>0\\
x+1>0
\end{array}\right.
\ \Longrightarrow \
\left\lbrace
\begin{array}{cc}
x>0\\
x>-1
\end{array}\right.
\ \Longrightarrow \
x\in (0,+\infty)
\\
\\
\left\lbrace
\begin{array}{cc}
x<0\\
x+1<0
\end{array}\right.
\ \Longrightarrow \
\left\lbrace
\begin{array}{cc}
x<0\\
x<-1
\end{array}\right.
\ \Longrightarrow \
x\in(-\infty,-1)
\end{array}\right\rbrace
\ \Longrightarrow \
x\in(-\infty,-1)\cup(0,+\infty)
\end{equation*}
Per tant la segona solució és $(-\infty,-1)\cup(0,+\infty)$.
\\
\\
\\
La solució de la inequació és la intersseció entre les dues solucions, és a dir:
\begin{equation*}
(-1,2)\cap((-\infty,-1)\cup(0,+\infty))=(0,2) \ \Longrightarrow \ x\in(0,2)
\end{equation*}